% Дүгнэлт

\phantomsection
\addchaptertocentry{Дүгнэлт}
\label{Conclusion}
\section*{\centering Дүгнэлт}

Энэхүү дипломын ажлаар рентген зургаас ясны хугарлыг автоматаар илрүүлэх, сегментлэх, хугарлын зэрэг тодорхойлох, эдгэрэлтийг таамаглах гэсэн дөрвөн даалгаврыг нэгтгэсэн олон даалгаварт (multi-task) гүн сургалтын систем боловсруулж, туршилтаар үнэлэв. EfficientNet-B3 хуваалцсан кодлогч дээр CoAtNet ангиллын модуль, U-Net сегментацийн модуль, ordinal regression хугарлын зэргийн модуль, Gaussian NLL regression эдгэрэлтийн модулийг каскад байдлаар холбосон Архитектур \#3-ийг хэрэгжүүлсэн бөгөөд нийт 13.2 сая параметртэй, NVIDIA RTX A6000 (48 GB) GPU-тай dedicated server дээр сургасан бөгөөд үнэлгээг Google Colab орчинд гүйцэтгэв.

Туршилтын үр дүнд ангилалын даалгаварт AUC-ROC = 0.8173, accuracy = 74.8\%, F1 = 0.734 буюу хүлээн авч болох гүйцэтгэл үзүүлсэн. Хугарлын зэрэглэлийн даалгаварт adjacent accuracy = 91.0\%, MAE = 0.535 буюу ordinal regression загвар клиникийн хугарлын зэрэглэлийг дарааллын хувьд алдаагүй таних чадварыг нотолсон. Эдгэрэлтийн модуль синтетик өгөгдөл дээр $R^2 = 0.46$ үзүүлж, хэсэгчлэн биелсэн үр дүн гарсан. Харин сегментацийн модуль нэгдсэн pipeline-д mean Dice = 0.0498 буюу системтэй бүтэлгүйтсэн бөгөөд үүний шалтгааныг маскийн талбайн хэт жижиг харьцаа ($\sim$0.2--0.4\%), background-давамгай loss-ийн локал минимум, $224\times224$ нягтрал дээр хугарлын нарийн шугам feature map-д алга болох зэрэг уялдсан хүчин зүйлсээр тайлбарлав. SAM zero-shot загварын Dice = 0.6427 болон тусдаа сургасан загварын Dice = 0.9633 нь capacity-ийн биш, сургалтын нэгтгэлийн асуудал гэдгийг баталгаажуулсан тул цаашдын судалгаанд focal Dice / Tversky loss болон өндөр нягтрал ашиглан шийдвэрлэх боломжтой.

Дипломын ажлын дөрвөн зорилтын хувьд: хугарал илрүүлэх ангилал болон хугарлын зэрэг тодорхойлох зорилт бүрэн биелсэн; эдгэрэлт таамаглах зорилт хэсэгчлэн биелсэн; сегментацийн зорилт бүрэн биелээгүй хэвээр үлдсэн. Эдгээр үр дүнг шударгаар хүлээн зөвшөөрч, хязгаарлалтуудыг тодорхой тогтоов: сегментацийн pipeline-ийн бүтэлгүйтэл цаашид заавал шийдвэрлэгдэх ёстой; эдгэрэлтийн модуль синтетик өгөгдөл дээр сургагдсан тул бодит клиникт баталгаажаагүй; Монгол хүн амын рентген зураг дээр гадны валидаци (external validation) хийгдээгүй байна.

Энэхүү судалгаа нь Монгол Улсад эмнэлгийн AI оношилгооны чиглэлийн анхдагч академик ажлуудын нэг болж байна. Радиологичдын тоо хязгаарлагдмал, хөдөө орон нутагт мэргэжилтний хүртээмж дутагдалтай орчинд эмчийн хяналтан дор ажиллах клиникийн шийдвэрийг дэмжих систем (CDS) нь сум, аймгийн эмнэлгийн зургийг урьдчилан шалгаж яаралтай тохиолдлыг эрэмбэлэх, залуу эмч оюутнуудын мэргэжлийн сургалтыг дэмжих боломж олгоно. Гэсэн хэдий ч recall = 72.7\% буюу хугарлын 27\% алдагдах нь өвчтөний аюулгүй байдлын үүднээс системийг бие даасан оношлогооны хэрэгсэл болгон хэрэглэхийг хориглох ба эмчийн эцсийн хяналт заавал шаардагдана.

NVIDIA RTX A6000 GPU дээр боловсруулсан үр ашигтай архитектур, Grad-CAM болон SHAP-аар тайлбарлах чадвартай загвар, болон шударгаар хийсэн бүтэлгүйтлийн дүн шинжилгээ нь дараагийн судлаачдад тодорхой замын зураглал үлдээж байна. Цаашдын хөгжүүлэлтэд сегментацийн модулийг яаралтай сайжруулах, Монголын бодит рентген зургаар гадаад баталгаажуулалт хийх, синтетик өгөгдлийг бодит эмнэлзүйн мэдээлэлд суурилсан өгөгдлөөр солих, CT болон MRI-г хамруулсан олон хэлбэрт (multi-modal) систем болгон өргөжүүлэх нь гол тэргүүлэх чиглэлүүд болно. Энэхүү судалгаа нь Монгол Улсын эрүүл мэндийн салбарт хиймэл оюун ухааны хэрэглээг ёс зүйтэй, хариуцлагатай, эмчийн хяналттай нэвтрүүлэх замын анхны алхам болж байгааг чухалчлан тэмдэглэнэ.
б